\begin{abstract}
The algorithm of Harrow, Hassidim and Lloyd (HHL) prepares a quantum state proportional to the
solution of a linear system $Ax=b$; it is the prototype of quantum linear algebra. We verify HHL in
Isabelle/HOL from end to end: from the spectral theorem, through phase estimation and the error
analysis of the original paper, to programs of quantum Hoare logic and circuits of one- and two-qubit
gates. In the exact case we prove total correctness of the repeat-until-success program: it
terminates with probability one in the normalised solution, after $1/p$ rounds on average. For
arbitrary spectra we prove the paper's Theorem~1 with explicit constants, and for the shifted cycle
Laplacian we compile the whole algorithm, state preparation and Hamiltonian simulation included, to
elementary gates with a closed-form gate count. Mechanising the paper's appendix refuted seven of its
statements, among them the constant of its Lemma~3, which must be $\pi^2$ rather than $\pi^2/2$.

The development has 130 theories and 64{,}000 lines, and an AI coding agent wrote almost all of it,
directed by one mathematician over 34 active days in six months, with a pause of 147 days in the
middle. HHL itself took eleven days, on top of a formalisation of a lecture course on quantum
computing. We report how the work was organised, what the human did and what the agent did, how the
agent failed, and what we learned about HHL and about formalising with an agent.
\keywords{Isabelle/HOL \and HHL algorithm \and Quantum Hoare logic \and Phase estimation \and
AI-assisted formalisation \and Formal verification}
\end{abstract}
